\documentclass[11pt,a4paper,sans]{moderncv}

\moderncvstyle{classic}
\usepackage{multicol}
\usepackage{lipsum}
\usepackage[hyperref]{}
\usepackage[scale=0.8, top=.7cm, bottom=.7cm, left=.8cm, right=.8cm]{geometry}

% makes a https hyperlink
% usage: \httpslink[optional text]{link}
\newcommand*{\httpslink}[2][]{% <=======================================
  \ifthenelse{\equal{#1}{}}%
    {\href{https://#2}{#2}}%
    {\href{https://#2}{#1}}}
    
\newcommand*{\googlescholarsocialsymbol}  {\includegraphics[width=0.25cm]{google-scholar}~} % <===================
\newcommand*{\researchgatesocialsymbol}  {\includegraphics[width=0.5cm]{example-image-b}~} % <===================
\newcommand*{\testsocialsymbol}  {\includegraphics[width=0.5cm]{example-image-c}~} % <===================

\makeatletter
\RenewDocumentCommand{\social}{O{}O{}m}{%
  \ifthenelse{\equal{#2}{}}%
    { 
    \ifthenelse{\equal{#1}{linkedin}}{\collectionadd[linkedin]{socials}{\protect\httpslink[#3]{www.linkedin.com/in/#3}}}{}%
    \ifthenelse{\equal{#1}{googlescholar}}{\collectionadd[googlescholar]{socials}{\protect\httpslink[#3]{scholar.google.com/citations?user=#3}}}{}% <================================================================
    \ifthenelse{\equal{#1}{researchgate}}{\collectionadd[researchgate]{socials}{\protect\httpslink[#3]{www.researchgate.com/profile/#3}}}{}% <================================================================
      \ifthenelse{\equal{#1}{twitter}} {\collectionadd[twitter]{socials} {\protect\httpslink[#3]{www.twitter.com/#3}}}    {}%
      \ifthenelse{\equal{#1}{github}}  {\collectionadd[github]{socials}  {\protect\httpslink[#3]{www.github.com/#3}}}     {}%
    }
    {\collectionadd[#1]{socials}{\protect\httpslink[#3]{#2}}}}
\makeatother


\usepackage[
    backend=biber,
    %defernumbers=true,
    sorting=ydnt]{biblatex}
\addbibresource{../../ags/meta/ags.bib}



% \defbibenvironment{bibliography} 
%   {\list
%      {\printtext[labelnumberwidth]{% label format from numeric.bbx
%         \printfield{labelprefix}%
%         \printfield{labelnumber}}}
%      {\setlength{\topsep}{0pt}% layout parameters from moderncvstyleclassic.sty
%       \setlength{\labelwidth}{\hintscolumnwidth}%
%       \setlength{\labelsep}{\separatorcolumnwidth}%
%       \leftmargin\labelwidth%
%       \advance\leftmargin\labelsep}%
%       \sloppy\clubpenalty4000\widowpenalty4000}
%   {\endlist}
%   {\item}

\moderncvstyle{classic}
\moderncvcolor{black}
\setlength{\hintscolumnwidth}{3.5cm} 
\setlength{\makecvheadnamewidth}{10cm}

\definecolor{linkColor}{rgb}{0,0,0}
\newcommand{\itlink}[2]{\textcolor{linkColor}{\textit{\link[#1]{#2}}}}
\newcommand{\bfitlink}[2]{\textcolor{linkColor}{\textbf{\textit{\link[#1]{#2}}}}}
\newcommand{\hpref}[2]{\hyperref[#1]{\textcolor{linkColor}{\textit{#2}}}}

\firstname{Aleksei G.}
\familyname{Sorokin}
\address{3443 S Morgan St, Chicago IL, 60608, US}
\phone[mobile]{+1~(630)~297~6261}  
\homepage{alegresor.github.io}
\email{asorokin@hawk.iit.edu}
\social[linkedin]{aleksei-sorokin}
\social[github]{alegresor}
\social[googlescholar]{akk3XSEAAAAJ}

\begin{document}
\makecvtitle
\vspace{-1.3cm}

\section{Skills}
\cvitem{\textbf{Math Interests}}{Probabilistic Numerics, Quasi-Monte Carlo, Gaussian Processes, Machine Learning}
\cvitem{\textbf{Softwares}}{GitHub Actions, Docker, Amazon EC2, \LaTeX}
\cvitem{\textbf{Languages}}{Python, Julia, C, MATLAB, R}
\cvitem{\textbf{Python Libraries}}{numpy, scipy, matplotlib, pandas, pytorch, gpytorch}

\section{Education}
\cvitem{2021 - 2025}{\textbf{Ph.D. in Applied Mathematics.} Illinois Institute of Technology (IIT).}
\cvitem{2017 - 2021}{\textbf{Master of Data Science.} IIT. Summa cum laude.}
\cvitem{2017 - 2021}{\textbf{B.S. in Applied Mathematics, Minor in Computer Science.} IIT. Summa cum laude.}

\section{Experiences}
\cvitem{Summer 2023}{\textbf{Graduate Intern} at \textbf{Los Alamos National Laboratory.} I modeled the solution processes of PDEs with random coefficients using efficient and error aware Gaussian processes. Resulted in publication of \citetitle{sorokin.gp4darcy}.}
\cvitem{Summer 2022}{\textbf{Givens Associate} at \textbf{Argonne National Laboratory}. I researched methods to efficiently estimate failure probability using Monte Carlo with non-parametric importance sampling. Resulted in publication of \citetitle{sorokin.adaptive_prob_failure_GP}.}
\cvitem{Summer 2021}{\textbf{Machine Learning Intern} at \textbf{SigOpt, an Intel Company}. I developed novel meta-learning techniques for model-aware hyperparameter optimization. Resulted in publication of \citetitle{sorokin.sigopt_mulch}.}
\cvitem{Fall 2021 - Present}{\textbf{Teaching assistant} at \textbf{IIT}. Includes leading review sessions for Ph.D. qualifying exams.}

\section{Projects}
\cvitem{\textbf{Fast Gaussian Processes with Derivatives}}{The cost of Gaussian process regression can be reduced from $\mathcal{O}(n^3)$ to $\mathcal{O}(n \log n)$ when one has control over the design of experiments. This is achieved by pairing quasi-random sampling sequences with matching kernels to induce structure in the kernel matrix. My Ph.D. research studies generalizations for quickly incorporating derivative information into the Gaussian process model and applications in building surrogates for PDEs solutions.}
\cvitem{\textbf{QMCPy}}{I have led development of QMCPy, a Quasi-Monte Carlo Python Library, since 2018. This package provides researchers and practitioners with high quality sequence generators, automatic variable transformations, adaptive stopping criteria algorithms, and diverse use cases. Publications include \citetitle{choi.QMC_software}, \citetitle{sorokin.MC_vector_functions_integrals}, \citetitle{choi.challenges_great_qmc_software}. See \itlink{qmcpy.org}{https://qmcpy.org} for details.}
\cvitem{\textbf{Argonne: AI on Supercomputers}}{I participated in an elective course on  \emph{AI Driven Science on Supercomputers} through \emph{Argonne National Laboratory}. Included instruction on handling large scale data pipelines and parallel training for neural networks. Coursework at \itlink{github.com/alegresor/ai-science-training-series}{https://github.com/alegresor/ai-science-training-series}.}

% \nocite{
%     sorokin.gp4darcy,
%     sorokin.MC_vector_functions_integrals,
%     choi.challenges_great_qmc_software,
%     sorokin.adaptive_prob_failure_GP,
%     sorokin.sigopt_mulch,
%     choi.QMC_software,
%     }
% \printbibliography[title=Publications]
    
%\bibliography{../../ags/meta/ags}

\section{Awards}
\cvitem{2023}{\textbf{Outstanding Mathematics Poster}, Los Alamos National Laboratory.}
\cvitem{2021}{\textbf{Best Manuscript}, IIT Undergraduate Research Journal.}
\cvitem{2020}{\textbf{Karl Menger Student Award for Exceptional Scholarship}, IIT.}

% \section{References}
% \cvitem{\href{mailto:hickernell@iit.edu}{hickernell@iit.edu}}{\textbf{Fred J. Hickernell, Ph.D.} Vice Provost for Research and Professor of Applied Mathematics, IIT.}
% \cvitem{\href{mailto:nickh@lanl.gov}{nickh@lanl.gov}}{\textbf{Nicolas W. Hengartner, Ph.D.} Senior Scientist, Los Alamos National Laboratory.}
% \cvitem{\href{mailto:mikemccourt1234@gmail.com}{mikemccourt1234@gmail.com}}{\textbf{Michael J. McCourt, Ph.D.} Co-Founder and CTO at Distributional.}
% \cvitem{\href{mailto:vhebbur@anl.gov}{vhebbur@anl.gov}}{\textbf{Vishwas Rao, Ph.D.} Assistant Computational Mathematician, Argonne National Laboratory.}

%%%%% COVER LETTER
%\clearpage
%\recipient{HR Department}{Corporation\\123 Pleasant Lane\\12345 City, State} % Letter recipient
%\date{\today} % Letter date
%\opening{Dear Sir or Madam,} % Opening greeting
%\closing{Sincerely yours,} % Closing phrase
%\enclosure[Attached]{curriculum vit\ae{}} % List of enclosed documents
%\makelettertitle % Print letter title
%\lipsum[1-3] % Dummy text
%\makeletterclosing % Print letter signature

\end{document}